\section{Naming a crate, and naming a device}
\label{sec:naming-and-device-selection}

Three names meet in a published crate, and confusing them is how a page ends up telling a reader to install something
that does not exist.

\begin{description}
  \item[The package name] is the identity on the registry: \texttt{efr32mg2x-pac}. It is what
        \texttt{cargo add} takes, what the URL uses, and what cannot be changed later without publishing again.
  \item[The library name] is what \texttt{use} reads: \texttt{efr32mg2x\_pac}. By default it is the package name with
        hyphens turned into underscores, which is exactly what a reader expects.
  \item[The device] is which part's register map you are compiling against. That belongs in a \emph{feature}, not in
        the crate name.
\end{description}

\subsection{Why the library name is allowed to differ, and why it should not}

Cargo resolves the name in a \texttt{use} statement from the library target, not from the package. When the dependency
entry is written with the package's own name, no rename happens and the library name wins; when it is written with a
different key, that key wins. So a crate called \texttt{efr32mg2x-pac} can export a library called
\texttt{efr32mg24\_pac}, and the following is correct:

\begin{lstlisting}[language=toml]
[dependencies]
efr32mg2x-pac = "0.1.3"        # the package
\end{lstlisting}

\begin{lstlisting}[language=Rust]
use efr32mg24_pac as pac;      # the library it exports
\end{lstlisting}

It compiles. It is also a trap: two different words for the same dependency in the same block of a README, and every
consumer has to learn which is which. This project did it deliberately for one release --- the package was renamed to
the family and the library kept the old name so that no dependent's \texttt{use} broke --- and then undid it, because the
moment to make the two agree is while every published version is yanked and no consumer exists yet.

The rule that follows: \textbf{make them the same}, and if you must differ, say so on the first screen of the README,
not in a comment halfway down.

\subsection{The device belongs in a feature}

A register map without a device is not a register map. The crate is generated for one part, the
EFR32MG24B220F1536IM48, and the feature names it:

\begin{lstlisting}[language=toml]
[features]
mg24 = []
default = ["mg24", "critical-section"]
\end{lstlisting}

Two things make that feature real rather than decorative. It is in \texttt{default}, so the ordinary dependency line
needs no ceremony. And its absence is an error, not a silent default:

\begin{lstlisting}[language=Rust]
#[cfg(not(feature = "mg24"))]
compile_error!("enable a device feature: `mg24` (EFR32MG24B220F1536IM48)");
\end{lstlisting}

A feature that selects nothing is worse than no feature: it promises a choice and delivers the same binary. The test for
any feature is whether turning it off changes something a reader can observe --- here, whether the crate compiles at all.

A second part then adds its own SVD, its own register modules behind its own feature, and the crate grows instead of
forking. Until that part exists, a crate that advertises twenty variants is advertising twenty things it cannot do; this
project shipped exactly that list, and every name in it was unusable.

\begin{keytakeaways}
  \item Package name, library name and device are three different things; the first two should read the same, and the
        third is a feature.
  \item \texttt{use} resolves to the library target. Understanding that is what makes the two-name trick work, and also
        what makes it a trap worth refusing.
  \item A device feature must change the build, or fail the build. Anything else is a promise with no code behind it.
  \item Rename while the registry has nothing usable of yours: yanked versions cost nothing to break, published ones
        cost a new version and a paragraph of explanation.
\end{keytakeaways}
